\chapter{Elementary inequalities}
\label{chap:pre_analysis}

This chapter collects the real-analysis facts used by the analysis of Entropic-BPI
(Chapters~\ref{chap:analysis_pos}, \ref{chap:analysis_neg} and~\ref{chap:sample_complexity}):
two square-root inequalities used to split the concentration terms, the bound
$(1 + a/H)^h \le e^a$ that removes the factor $\kappa^{h-1}$ of the unrolled certificate
recursion, the Bhatia--Davis inequality and the normalized variance bound for exponential
transforms (the paper's Lemma~24, \cite{essakine2026tight}), the logarithmic sum bound (the
paper's Lemma~28) used by the counting argument, and the self-bounding inequality (the paper's
Lemma~29, Lemma~13 of \cite{menard2021fast}) that turns the recursive bound on the number of
episodes into an explicit one. Everything is stated over $\R$ with no reference to the MDP.

\textbf{Constants.} Every constant is verified below. Two of the paper's statements are
adjusted: the logarithmic sum bound (Lemma~\ref{lem:log_sum_bound}) holds with the constant
$3$ in place of $4$ (both forms are stated, since the form with $4$ is the one used by the
counting argument); and the self-bounding inequality
(Lemma~\ref{lem:self_bounding}) is \emph{false} with the constant $C_1$ printed in the paper
(a counterexample is given in Remark~\ref{rem:pana_self_bounding_paper}), so it is stated and
proved with a corrected $C_1$.

\textbf{What Mathlib provides.} The real square root \texttt{Real.sqrt} with
\texttt{Real.sqrt\_le\_sqrt}, \texttt{Real.sqrt\_le\_left}, \texttt{Real.le\_sqrt},
\texttt{Real.sq\_sqrt}, \texttt{Real.sqrt\_mul}; the exponential and logarithm with
\texttt{Real.add\_one\_le\_exp}, \texttt{Real.log\_le\_sub\_one\_of\_pos},
\texttt{Real.one\_sub\_inv\_le\_log\_of\_pos}, \texttt{Real.log\_le\_log\_iff},
\texttt{Real.log\_rpow}, \texttt{Real.rpow\_natCast}; the variance
\texttt{ProbabilityTheory.variance} with \texttt{variance\_eq\_sub} (variance as
$\E X^2 - (\E X)^2$ for $X \in L^2$), \texttt{memLp\_of\_bounded} and the Bhatia--Davis
inequality itself (\texttt{ProbabilityTheory.variance\_le\_sub\_mul\_sub},
Lemma~\ref{lem:bhatia_davis}); the AM--GM inequality in the form \texttt{two\_mul\_le\_add\_sq};
the weak logarithm-versus-power bound \texttt{Real.log\_le\_rpow\_div}.

\textbf{What is missing.} The normalized variance bound, the logarithmic sum bound and the
self-bounding inequality, and the small lemmas around them. They are in
\texttt{Essakine2026Tight/Mathlib/} (\texttt{Algebra/Order/Field/Basic.lean},
\texttt{Analysis/Real/Sqrt.lean}, \texttt{Analysis/SpecialFunctions/Exp.lean},
\texttt{Analysis/SpecialFunctions/Log/Basic.lean}, \texttt{Analysis/SpecialFunctions/Pow/Real.lean},
\texttt{Analysis/SpecialFunctions/Log/SelfBounding.lean},
\texttt{Probability/Moments/Variance.lean}).

\section{Square roots, exponentials and logarithms}
\label{sec:pre_analysis_basic}

\begin{lemma}[Subadditivity of the square root]\label{lem:sqrt_add_le}
\lean{Real.sqrt_add_le}\leanok
For $a, b \ge 0$, $\sqrt{a + b} \le \sqrt a + \sqrt b$.
\end{lemma}

\begin{proof}
\leanok
Both sides are nonnegative, so it suffices to compare squares (Mathlib
\texttt{Real.sqrt\_le\_left} with $y = \sqrt a + \sqrt b \ge 0$):
$(\sqrt a + \sqrt b)^2 = a + b + 2\sqrt a\sqrt b \ge a + b$ because $\sqrt a\sqrt b \ge 0$
(\texttt{Real.sq\_sqrt}, \texttt{Real.sqrt\_nonneg}).
\end{proof}

\begin{lemma}[Weighted AM--GM for a square root]\label{lem:sqrt_mul_le}
\lean{Real.two_mul_sqrt_mul_le, Real.sqrt_mul_le_half_add_half, Real.sqrt_mul_le_add}\leanok
For $a, b \ge 0$ and $\eta > 0$,
\[
  \sqrt{ab} \le \frac{\eta a}{2} + \frac{b}{2\eta} \le \eta a + \frac{b}{\eta} .
\]
\end{lemma}

\begin{proof}
\leanok
Let $x := \sqrt{\eta a}$ and $y := \sqrt{b/\eta}$, so that $x, y \ge 0$, $x^2 = \eta a$,
$y^2 = b/\eta$ and $xy = \sqrt{ab}$ (\texttt{Real.sqrt\_mul}, \texttt{Real.sq\_sqrt}; $\eta a \ge 0$
and $b/\eta \ge 0$). From $(x - y)^2 \ge 0$ we get $2xy \le x^2 + y^2$
(\texttt{two\_mul\_le\_add\_sq}), i.e.\ $2\sqrt{ab} \le \eta a + b/\eta$
(\texttt{Real.two\_mul\_sqrt\_mul\_le}), which is the first
inequality. The second holds because $\eta a \ge 0$ and $b/\eta \ge 0$.
\end{proof}

\begin{lemma}[A power bound for $(1 + a/H)^h$]\label{lem:one_add_div_pow_le_exp}
\lean{Real.one_add_div_pow_le_exp}\leanok
For $a \ge 0$, integers $H \ge 1$ and $0 \le h \le H$: $(1 + a/H)^h \le e^{a}$.
\end{lemma}

\begin{proof}
\leanok
By \texttt{Real.add\_one\_le\_exp}, $1 + a/H \le e^{a/H}$, and $1 + a/H \ge 0$. Since
$x \mapsto x^h$ is monotone on $[0, \infty)$ (\texttt{pow\_le\_pow\_left₀}),
$(1 + a/H)^h \le (e^{a/H})^h = e^{ah/H}$ (\texttt{Real.exp\_nat\_mul}), and $ah/H \le a$
because $0 \le h \le H$ and $a \ge 0$, so $e^{ah/H} \le e^{a}$ (\texttt{Real.exp\_le\_exp}).
\end{proof}

\begin{lemma}[{$2 - x \le 1/x$ on $(0, \infty)$}]\label{lem:two_sub_le_inv}
\lean{two_sub_le_inv}\leanok
For $x > 0$ (in any linearly ordered field), $2 - x \le 1/x$. (It is used for
$x = e^{\beta\eps} \in (0, 1]$.)
\end{lemma}

\begin{proof}
\leanok
$1/x - (2 - x) = (1 - x)^2/x \ge 0$ (\texttt{field\_simp}, \texttt{positivity}).
\end{proof}

\section{The Bhatia--Davis inequality and the normalized variance bound}
\label{sec:pre_analysis_variance}

\begin{lemma}[Bhatia--Davis inequality]\label{lem:bhatia_davis}
\lean{ProbabilityTheory.variance_le_sub_mul_sub}\leanok
Let $(\Omega, \Prob)$ be a probability space, $Y : \Omega \to \R$ an a.e.-measurable random
variable and $m, M$ real numbers with $Y \in [m, M]$ almost surely. Then, with $\mu := \E Y$,
\[
  \Var(Y) \le (M - \mu)(\mu - m) .
\]
\end{lemma}

\begin{proof}
\leanok
This is Mathlib's \texttt{ProbabilityTheory.variance\_le\_sub\_mul\_sub}. The argument:
$Y$ is bounded a.s.\ on a probability space, hence in $L^2$
(\texttt{memLp\_of\_bounded}), and $\Var(Y) = \E Y^2 - \mu^2$ (\texttt{variance\_eq\_sub}).
Almost surely $(M - Y)(Y - m) \ge 0$, so
\[
  0 \le \E[(M - Y)(Y - m)] = \E[-Y^2 + (M + m)Y - Mm] = -\E Y^2 + (M + m)\mu - Mm ,
\]
i.e.\ $\E Y^2 \le (M + m)\mu - Mm$. Therefore
$\Var(Y) = \E Y^2 - \mu^2 \le (M + m)\mu - Mm - \mu^2 = (M - \mu)(\mu - m)$.
This is the inequality of \cite{bhatia2000better}.
\end{proof}

\begin{lemma}[The mean of a bounded random variable]\label{lem:pana_integral_mem_Icc}
\lean{ProbabilityTheory.integral_mem_Icc_of_ae_mem_Icc}\leanok
Let $(\Omega, \Prob)$ be a probability space and $Y : \Omega \to \R$ a.e.-strongly measurable
with $Y \in [m, M]$ almost surely. Then $\E Y \in [m, M]$.
\end{lemma}

\begin{proof}
\leanok
$Y$ is integrable (\texttt{memLp\_of\_bounded} with $p = 1$), and the integral is monotone
(\texttt{integral\_mono\_ae}) with $\E[m] = m$, $\E[M] = M$ on a probability space.
\end{proof}

\textbf{Lean remark.} The finite-vector version $\Var_p(f) \le (M - pf)(pf - m)$ for a
probability vector $p$ on a finite set and $f \in [m, M]$ is the same statement for the measure
$\sum_s p(s)\delta_s$ (\texttt{weightedMeasure}), or is proved directly by the same computation
with finite sums.

\begin{lemma}[Maximization of the normalized Bhatia--Davis bound]\label{lem:pana_var_ratio_max}
\lean{sub_mul_sub_div_sq_le}\leanok
For $m, M, \mu > 0$ (in any linearly ordered field),
\[
  \frac{(M - \mu)(\mu - m)}{\mu^2} \le \frac{(M - m)^2}{4Mm}
  = \frac{(M/m - 1)^2}{4\,M/m},
\]
with equality if and only if $\mu = 2Mm/(M + m)$.
\end{lemma}

\begin{proof}
\leanok
Since $\mu > 0$ and $4Mm > 0$, the inequality is equivalent to
$4Mm(M - \mu)(\mu - m) \le (M - m)^2\mu^2$ (\texttt{div\_le\_div\_iff₀}). Expanding both sides,
\[
  (M - m)^2\mu^2 - 4Mm(M - \mu)(\mu - m)
  = (M + m)^2\mu^2 - 4Mm(M + m)\mu + 4M^2m^2
  = \big((M + m)\mu - 2Mm\big)^2 \ge 0 ,
\]
where we used $(M - m)^2 + 4Mm = (M + m)^2$ (\texttt{nlinarith} with this square as hint).
Equality holds iff $(M + m)\mu = 2Mm$. The last
expression of the statement is obtained by dividing numerator and denominator by $m^2$.
(This is the calculus maximization of the paper's proof of Lemma~24 made algebraic: the
function $\mu \mapsto (M - \mu)(\mu - m)/\mu^2$ increases on $[m, 2Mm/(M + m)]$ and decreases
on $[2Mm/(M + m), M]$, with maximum $(M - m)^2/(4Mm)$. The restriction $\mu \in [m, M]$ is not
needed.)
\end{proof}

\begin{lemma}[Normalized variance bound for an exponential transform]\label{lem:normalized_variance_bound}
\uses{lem:bhatia_davis, lem:pana_var_ratio_max, lem:pana_integral_mem_Icc}
\lean{ProbabilityTheory.variance_exp_div_sq_integral_le}\leanok
Let $(\Omega, \Prob)$ be a probability space, $R \ge 0$, $\beta \in \R$ and $f : \Omega \to \R$
an a.e.-measurable random variable with $f \in [0, R]$ almost surely. Let $Y := e^{\beta f}$
and $\mu := \E Y$. Then
\[
  \frac{\Var(Y)}{\mu^2} \le \frac{(e^{\abs\beta R} - 1)^2}{4\,e^{\abs\beta R}} .
\]
\end{lemma}

\begin{proof}
\uses{lem:bhatia_davis, lem:pana_var_ratio_max, lem:pana_integral_mem_Icc}
\leanok
Let $m := \min\{1, e^{\beta R}\}$ and $M := \max\{1, e^{\beta R}\}$.
Almost surely $\beta f$ lies between $0$ and $\beta R$ (if $\beta \ge 0$ then
$0 \le \beta f \le \beta R$; if $\beta < 0$ then $\beta R \le \beta f \le 0$), so by
monotonicity of $\exp$, $m \le Y \le M$ a.s.\ (\texttt{ProbabilityTheory.exp\_mul\_mem\_Icc\_min\_max}).
Lemma~\ref{lem:pana_integral_mem_Icc} gives
$\mu \in [m, M]$, and $m > 0$, hence $\mu > 0$; Lemma~\ref{lem:bhatia_davis} gives
$\Var(Y) \le (M - \mu)(\mu - m)$.
Dividing by $\mu^2$ and applying Lemma~\ref{lem:pana_var_ratio_max},
\[
  \frac{\Var(Y)}{\mu^2} \le \frac{(M - \mu)(\mu - m)}{\mu^2} \le \frac{(M - m)^2}{4Mm} .
\]
Finally $(M - m)^2/(4Mm) = (e^{\abs\beta R} - 1)^2/(4e^{\abs\beta R})$
(\texttt{ProbabilityTheory.max\_sub\_min\_sq\_div\_eq}): if $\beta \ge 0$ then $m = 1$,
$M = e^{\beta R}$; if $\beta < 0$ then $m = e^{\beta R}$, $M = 1$ and
$(1 - e^{\beta R})^2/(4e^{\beta R}) = (e^{-\beta R} - 1)^2/(4e^{-\beta R})$.
\end{proof}

\textbf{Lean remark.} This is the paper's Lemma~24 (Appendix~F). The bound is used in
Lemma~\ref{lem:variance_ratio_le} with $f$ the return of an episode and $R = \Gmax$; the
paper's constant $V_G = (e^{\abs\beta G} - 1)^2/e^{\abs\beta G}$ (Definition~\ref{def:upper_bound_const})
is four times the bound above, so the factor $4$ is slack that the formal proof may drop.
The case split on the sign of $\beta$ only appears in the two auxiliary lemmas (the a.s.\ bounds
on $Y$ and the identification of $(M - m)^2/(4Mm)$).

\section{The logarithmic sum bound}
\label{sec:pre_analysis_log_sum}

\begin{lemma}[One step of the logarithmic sum bound]\label{lem:pana_log_sum_step}
\lean{Real.div_max_one_le_three_mul_log}\leanok
For $U \ge 0$ and $u \in [0, 1]$,
\[
  \frac{u}{\max\{U, 1\}} \le 3\log\frac{U + u + 1}{U + 1} .
\]
\end{lemma}

\begin{proof}
\leanok
Let $x := u/(U + 1) \ge 0$. Since $1 + x > 0$, Mathlib \texttt{Real.one\_sub\_inv\_le\_log\_of\_pos}
gives $\log(1 + x) \ge 1 - \frac{1}{1 + x} = \frac{x}{1 + x} = \frac{u}{U + 1 + u}$, so the
right-hand side is at least $3u/(U + 1 + u)$. It
suffices to show $u(U + 1 + u) \le 3u\max\{U, 1\}$, which follows from $u \ge 0$ and
$U + 1 + u \le 3\max\{U, 1\}$: if $U \ge 1$ then
$U + 1 + u \le U + 2 \le 3U$; if $U < 1$ then $U + 1 + u < 3 = 3\max\{U, 1\}$. Hence
$\frac{u}{\max\{U,1\}} \le \frac{3u}{U + 1 + u} \le 3\log\frac{U + u + 1}{U + 1}$.
\end{proof}

\begin{lemma}[Logarithmic sum bound]\label{lem:log_sum_bound}
\uses{lem:pana_log_sum_step}
\lean{Real.sum_div_max_one_le_three_mul_log, Real.sum_div_max_one_le_four_mul_log}\leanok
Let $(u_t)_{t \ge 1}$ be a sequence with $u_t \in [0, 1]$ and $U_t := \sum_{i = 1}^t u_i$
($U_0 = 0$). For every integer $T \ge 0$,
\[
  \sum_{t = 0}^{T - 1}\frac{u_{t + 1}}{\max\{U_t, 1\}} \le 3\log(U_{T} + 1) \le 4\log(U_{T} + 1) .
\]
\end{lemma}

\begin{proof}
\uses{lem:pana_log_sum_step}
\leanok
For each $t$, $U_t \ge 0$, $u_{t+1} \in [0,1]$ and $U_{t+1} = U_t + u_{t+1}$, so
Lemma~\ref{lem:pana_log_sum_step} gives
\[
  \frac{u_{t + 1}}{\max\{U_t, 1\}} \le 3\log\frac{U_{t+1} + 1}{U_t + 1}
  = 3\big(\log(U_{t+1} + 1) - \log(U_t + 1)\big)
\]
(\texttt{Real.log\_div}, both arguments being positive). Summing over $t = 0, \dots, T - 1$, the
right-hand side telescopes (\texttt{Finset.sum\_range\_sub}) to
$3(\log(U_{T} + 1) - \log(U_0 + 1)) = 3\log(U_{T} + 1)$. Finally $U_{T} + 1 \ge 1$, so
$\log(U_{T} + 1) \ge 0$ (\texttt{Real.log\_nonneg}) and $3\log(\cdot) \le 4\log(\cdot)$.
\end{proof}

\begin{remark}[On the constant]\label{rem:pana_log_sum_const}
This is the paper's Lemma~28 (Appendix~F), stated there with the constant $4$ and without
proof (it is Lemma~8 of \cite{menard2021fast}). The one-step bound holds with the constant
$3$, hence so does the sum bound; the constant $2$ fails at $U = u = 1$ ($1 > 2\log(3/2)$).
The form with the constant $4$ is the one consumed by
Lemma~\ref{lem:sum_counts_bound}. \textbf{Lean remark.} The sequence is indexed by $\N$ with
$u_{t+1} = $ \texttt{u t}, so that $U_t = $ \texttt{∑ i ∈ Finset.range t, u i} and the sum is
\texttt{∑ t ∈ Finset.range T, u t / max (∑ i ∈ Finset.range t, u i) 1}.
\end{remark}

\section{The self-bounding inequality}
\label{sec:pre_analysis_self_bounding}

The counting argument of Chapter~\ref{chap:sample_complexity} produces an inequality of the
form $\tau \le C\sqrt{\tau\,\ell(\tau)} + D\,\ell'(\tau)$ with $\ell, \ell'$ polylogarithmic in
$\tau$, and one must solve it for $\tau$. We first solve the quadratic part, then bound the
logarithm of a quantity that is bounded by a multiple of the square of its own logarithm.

\begin{lemma}[Root of a quadratic inequality]\label{lem:pana_quadratic_root}
\lean{Real.le_add_sqrt_of_sq_le}\leanok
Let $b, c \ge 0$ and $x$ real with $x^2 \le bx + c$. Then $x \le b + \sqrt c$.
\end{lemma}

\begin{proof}
\leanok
Suppose $x > b + \sqrt c$. Then $x > 0$ (as $b + \sqrt c \ge 0$), $x > b$ and $x > \sqrt c \ge 0$,
so
\[
  x^2 = x\cdot x > x(b + \sqrt c) = bx + x\sqrt c \ge bx + \sqrt c\sqrt c = bx + c ,
\]
using $x\sqrt c \ge \sqrt c\sqrt c$ (from $x \ge \sqrt c$ and $\sqrt c \ge 0$) and
$\sqrt c\sqrt c = c$ (\texttt{Real.mul\_self\_sqrt}). This contradicts $x^2 \le bx + c$.
(No sign assumption on $x$ is needed.)
\end{proof}

\begin{lemma}[Logarithm versus powers]\label{lem:pana_log_le_rpow}
\lean{Real.log_le_rpow_div_exp_one_mul}\leanok
For $y > 0$ and $\gamma > 0$, $\log y \le \dfrac{y^{\gamma}}{e\gamma}$.
\end{lemma}

\begin{proof}
\leanok
Let $z := y^\gamma > 0$. Mathlib \texttt{Real.log\_le\_sub\_one\_of\_pos} applied to $z/e > 0$
gives $\log(z/e) \le z/e - 1$, and $\log(z/e) = \log z - 1$ (\texttt{Real.log\_div},
\texttt{Real.log\_exp}), so $\log z \le z/e$. Since $\log z = \gamma\log y$
(\texttt{Real.log\_rpow}), $\gamma\log y \le y^\gamma/e$; divide by $\gamma > 0$.
(Mathlib's \texttt{Real.log\_le\_rpow\_div} is the same bound without the factor $1/e$.)
\end{proof}

\begin{lemma}[A quantity bounded by the square of its logarithm]\label{lem:pana_log_self_bound}
\uses{lem:pana_log_le_rpow}
\lean{Real.log_le_of_le_mul_sq_log}\leanok
Let $y \ge 1$ and $M > 0$ with $y \le M(\log y)^2$. Then
\[
  \log y \le \tfrac85\log(4M) .
\]
\end{lemma}

\begin{proof}
\uses{lem:pana_log_le_rpow}
\leanok
Since $y \ge 1$, $\log y \ge 0$, and Lemma~\ref{lem:pana_log_le_rpow} with $\gamma = 3/16$ gives
$0 \le \log y \le \frac{16}{3e}\,y^{3/16}$. Squaring (both sides are nonnegative,
\texttt{pow\_le\_pow\_left₀}) and using $y^{3/16}\cdot y^{3/16} = y^{3/8}$ (\texttt{Real.rpow\_add}),
\[
  (\log y)^2 \le \frac{256}{9e^2}\,y^{3/8} \le 4\,y^{3/8},
\]
because $9e^2 \ge 64$ (indeed $e \ge 2.7$, \texttt{Real.exp\_one\_gt\_d9}, so $e^2 \ge 7.29 > 64/9$).
Hence $y \le 4M\,y^{3/8}$, and since $y = y^{5/8}\,y^{3/8}$ with $y^{3/8} > 0$
(\texttt{Real.rpow\_add}, \texttt{le\_of\_mul\_le\_mul\_right}), $y^{5/8} \le 4M$. Both sides are
positive, so taking logarithms (\texttt{Real.log\_le\_log}, \texttt{Real.log\_rpow}) yields
$\frac58\log y \le \log(4M)$, i.e.\ $\log y \le \frac85\log(4M)$.
\end{proof}

\begin{lemma}[Reduction of the self-bounding inequality to a quadratic one]\label{lem:pana_self_bounding_quadratic}
\uses{lem:pana_quadratic_root}
\lean{Real.le_mul_sq_log_of_le_sqrt_mul_add}\leanok
Let $A, B, C, D \ge 0$, $E \ge B$, $\alpha \in \R$ and $\tau \ge 0$ with $L := \log(\alpha\tau) \ge 1$
satisfy
\begin{equation}\label{eq:pana_self_bounding_hyp}
  \tau \le C\sqrt{\tau\big(A\log(\alpha\tau) + B\log(\alpha\tau)^2\big)}
  + D\big(A\log(\alpha\tau) + E\log(\alpha\tau)^2\big) .
\end{equation}
Let $K := C^2(A + B) + (D + 2\sqrt D\,C)(A + E)$. Then $\tau \le K L^2$.
\end{lemma}

\begin{proof}
\uses{lem:pana_quadratic_root}
\leanok
\begin{enumerate}
  \item[(i)] \emph{Reduction to a quadratic inequality.} Since $L \ge 1$ and $A \ge 0$,
        $AL \le AL^2$, hence $AL + BL^2 \le (A + B)L^2$ and $AL + EL^2 \le (A + E)L^2$. Both
        terms of~\eqref{eq:pana_self_bounding_hyp} are monotone in these quantities ($C, D \ge 0$,
        $\sqrt{\cdot}$ monotone), and $\sqrt{\tau(A + B)L^2} = L\sqrt{A+B}\sqrt\tau$ as $L \ge 0$
        (\texttt{Real.sqrt\_mul}, \texttt{Real.sqrt\_sq}); so
        $\tau \le CL\sqrt{A + B}\,\sqrt\tau + D(A + E)L^2$.
  \item[(ii)] \emph{Solving it.} Put $x := \sqrt\tau$, $b := CL\sqrt{A+B} \ge 0$ and
        $c := D(A+E)L^2 \ge 0$; then $x^2 = \tau \le bx + c$ and Lemma~\ref{lem:pana_quadratic_root}
        gives $x \le b + \sqrt c = CL\sqrt{A+B} + \sqrt D\sqrt{A+E}\,L$. Squaring (both sides
        nonnegative),
        \[
          \tau \le L^2\Big(C^2(A+B) + D(A+E)\Big) + 2C\sqrt D\,L^2\,\sqrt{A+B}\sqrt{A+E} .
        \]
        As $B \le E$, $\sqrt{A+B} \le \sqrt{A+E}$, so $\sqrt{A+B}\sqrt{A+E} \le A + E$ and
        the last term is at most $2C\sqrt D\,(A+E)L^2$. Hence $\tau \le KL^2$.
\end{enumerate}
\end{proof}

\begin{lemma}[The logarithm of the self-bounding inequality]\label{lem:pana_self_bounding_log}
\uses{lem:pana_self_bounding_quadratic, lem:pana_log_self_bound}
\lean{Real.one_le_log_mul_of_exp_one_le, Real.log_le_of_le_sqrt_mul_add}\leanok
Let $A, B, C, D \ge 0$, $E \ge B$, $\alpha \ge e$ and $\tau > 1$
satisfy~\eqref{eq:pana_self_bounding_hyp}, and let $K$ be as in
Lemma~\ref{lem:pana_self_bounding_quadratic}. Then $L := \log(\alpha\tau) \ge 1$, $K > 0$ and
$L \le \frac85\log(4\alpha K)$.
\end{lemma}

\begin{proof}
\uses{lem:pana_self_bounding_quadratic, lem:pana_log_self_bound}
\leanok
$\alpha\tau \ge \alpha \ge e$, so $L \ge 1$ (\texttt{Real.le\_log\_iff\_exp\_le};
\texttt{Real.one\_le\_log\_mul\_of\_exp\_one\_le}). Lemma~\ref{lem:pana_self_bounding_quadratic}
gives $\tau \le KL^2$; since $\tau > 0$, this forces $K > 0$. Let $y := \alpha\tau \ge 1$ and
$M := \alpha K > 0$. Multiplying $\tau \le KL^2$ by $\alpha$ gives $y \le M(\log y)^2$, and
Lemma~\ref{lem:pana_log_self_bound} yields $L = \log y \le \frac85\log(4M) = \frac85\log(4\alpha K)$.
\end{proof}

\begin{lemma}[Self-bounding inequality, sharp form]\label{lem:pana_self_bounding_core}
\uses{lem:pana_self_bounding_quadratic, lem:pana_self_bounding_log}
\lean{Real.le_mul_sq_log_add_one_of_le_sqrt_mul_add}\leanok
Let $A, B, C, D \ge 0$, $E \ge B$, $\alpha \ge e$ and $\tau \ge 0$
satisfy~\eqref{eq:pana_self_bounding_hyp}, let $K := C^2(A + B) + (D + 2\sqrt D\,C)(A + E)$ and
let $M \ge 4\alpha K$. Then
\[
  \tau \le K\Big(\tfrac85\log M\Big)^2 + 1 .
\]
\end{lemma}

\begin{proof}
\uses{lem:pana_self_bounding_quadratic, lem:pana_self_bounding_log}
\leanok
If $\tau \le 1$ the claim holds because $K \ge 0$. Assume $\tau > 1$. By
Lemma~\ref{lem:pana_self_bounding_log}, $K > 0$, $L := \log(\alpha\tau) \ge 1$ and
$L \le \frac85\log(4\alpha K) \le \frac85\log M$ ($0 < 4\alpha K \le M$ and $\log$ is monotone);
by Lemma~\ref{lem:pana_self_bounding_quadratic}, $\tau \le KL^2$.
Then $0 \le L \le \frac85\log M$ gives $L^2 \le (\frac85\log M)^2$, and
$\tau \le KL^2 \le K(\frac85\log M)^2 \le K(\frac85\log M)^2 + 1$.
\end{proof}

\begin{lemma}[The constant of the self-bounding inequality]\label{lem:pana_four_alpha_K_le}
\lean{Real.four_mul_mul_le}\leanok
Let $A, B, C, D \ge 0$, $E \ge B$ and $\alpha \ge e$, and $K := C^2(A + B) + (D + 2\sqrt D\,C)(A + E)$.
Then $4\alpha K \le 11\alpha^2(A + E)(C + \sqrt D)^2$.
\end{lemma}

\begin{proof}
\leanok
$K \le C^2(A+E) + (D + 2\sqrt D C)(A+E) = (A+E)\big(C^2 + 2C\sqrt D + D\big) = (A+E)(C + \sqrt D)^2$
(using $B \le E$ and $\sqrt D^2 = D$), and $4\alpha \le 11\alpha^2$ because $\alpha \ge e \ge 4/11$.
Both factors are nonnegative, so the products compare.
\end{proof}

\begin{lemma}[Self-bounding inequality]\label{lem:self_bounding}
\uses{lem:pana_self_bounding_core, lem:pana_four_alpha_K_le}
\lean{Real.le_of_le_sqrt_mul_add, Real.le_of_le_sqrt_mul_add'}\leanok
Let $A, B, C, D \ge 0$, $E \ge B$, $\alpha \ge e$ and $\tau \ge 0$ satisfy
\[
  \tau \le C\sqrt{\tau\big(A\log(\alpha\tau) + B\log(\alpha\tau)^2\big)}
  + D\big(A\log(\alpha\tau) + E\log(\alpha\tau)^2\big) .
\]
Then
\[
  \tau \le C^2(A + B)\,C_1^2 + \big(D + 2\sqrt D\,C\big)(A + E)\,C_1^2 + 1,
  \qquad
  C_1 := \frac85\log\Big(11\,\alpha^2\,(A + E)\,\big(C + \sqrt D\big)^2\Big) .
\]
If moreover $D \ge 1$, the same holds with $C_1 := \frac85\log\big(11\alpha^2(A+E)(C + D)^2\big)$.
\end{lemma}

\begin{proof}
\uses{lem:pana_self_bounding_core, lem:pana_four_alpha_K_le}
\leanok
Let $K := C^2(A+B) + (D + 2\sqrt D C)(A+E)$, so that the claimed bound reads
$\tau \le KC_1^2 + 1$, which is Lemma~\ref{lem:pana_self_bounding_core} with
$M := 11\alpha^2(A+E)(C+\sqrt D)^2 \ge 4\alpha K$ (Lemma~\ref{lem:pana_four_alpha_K_le}).
If $D \ge 1$ then $\sqrt D \le D$ (\texttt{Real.sqrt\_le\_left}), so
$(C + \sqrt D)^2 \le (C + D)^2$ and $M' := 11\alpha^2(A+E)(C+D)^2 \ge M \ge 4\alpha K$ also
qualifies.
\end{proof}

\begin{remark}[The paper's constant is wrong]\label{rem:pana_self_bounding_paper}
The paper's Lemma~29 (Appendix~F, quoted from Lemma~13 of \cite{menard2021fast}) states the
conclusion with $C_1 = \frac85\log(11\alpha^2(A+E)(C+D))$, i.e.\ with $(C + D)$ instead of
$(C + \sqrt D)^2$, for positive $A, B, C, D, E$ with $1 \le B \le E$. That statement is false.
Take $\alpha = e$, $B = E = 1$, $C = 1000$, $A = D = 10^{-6}$ and $\tau = 4\cdot 10^8$. Then
$L = \log(\alpha\tau) \approx 20.81$ and the right-hand side of the hypothesis is
$C\sqrt{\tau(AL + L^2)} + D(AL + L^2) \approx 1000\cdot\sqrt{4\cdot10^8}\cdot 20.81
\approx 4.16\cdot 10^8 \ge \tau$, so the hypothesis holds; but the paper's
$C_1 = \frac85\log(11e^2(1 + 10^{-6})(1000 + 10^{-6})) \approx 18.09$ gives the bound
$C^2(A+B)C_1^2 + (D + 2\sqrt D C)(A+E)C_1^2 + 1 \approx 10^6\cdot 327.2 + 655 \approx 3.27\cdot 10^8 < \tau$.
(The largest $\tau$ satisfying the hypothesis is about $4.37\cdot 10^8$; the corrected $C_1$
of Lemma~\ref{lem:self_bounding} is $\approx 29.1$ and gives the bound $\approx 8.5\cdot10^8$.)
The defect is in the exponent of $C + D$: for $D \to 0$ and $A \to 0$ the hypothesis
$\tau \le C\sqrt\tau\,L$ only gives $\tau \le C^2L^2$ with $L \approx 2\log C + \log\alpha +
2\log L$, which is not below $\frac85\log C + O(1)$ for large $C$. The correct exponent is
$2$, and this is what Lemma~\ref{lem:self_bounding} proves, with the weaker hypotheses
$A, B, C, D \ge 0$, $B \le E$ (positivity and $B \ge 1$ are not used). The sharper bound
$C_1 = \frac85\log(4\alpha K)$ of Lemma~\ref{lem:pana_self_bounding_core} is what the proof
actually gives. In the application (Definition~\ref{def:upper_bound_const}), $D \ge 1$, so
the form with $(C + D)^2$ applies; the asymptotic order of the bound is unchanged (only the
logarithmic factor changes).

\textbf{Lean remark.} \texttt{Real.log 0 = 0}, so the hypothesis and the conclusion make sense
for $\tau = 0$ and $K = 0$ without side conditions; the case $\tau \le 1$ is treated first and
$K > 0$ is derived in the other case. The powers $y^{3/16}$, $y^{3/8}$, $y^{5/8}$ are
\texttt{Real.rpow}; $\sqrt\tau$ is \texttt{Real.sqrt} with \texttt{Real.sq\_sqrt},
\texttt{Real.sqrt\_le\_sqrt}. The numerical facts are $e \ge 2.7$
(\texttt{Real.exp\_one\_gt\_d9}) and $4/11 \le e$.
\end{remark}
